\chapter{Methodology}
\label{ch:methodology}

This chapter describes the data pipeline, the model, the procedure by which the
operating configuration was selected, and the implementation of the Spike
Activation Map. Every measurement reported here is on the validation partition.

\section{Data}

\subsection{Acquisition}

The Gravity Spy machine-learning release~\cite{glanzer2021} is distributed as
eight CSV files, one per detector (H1, L1) and observing run (O1, O2, O3a, O3b).
Each record corresponds to an Omicron trigger with signal-to-noise ratio above
$7.5$ and peak frequency between $10$\,Hz and $2048$\,Hz, and carries the
trigger metadata, the per-class confidence of the Gravity Spy classifier, the
assigned label, and URLs to Omega-scan spectrograms.

\textbf{The labels are the classifications produced by that convolutional
model}, not physical ground truth and not the volunteer-combined label of the
separate citizen-science release. Accuracy reported in this thesis is agreement
with that reference classifier, which is standard practice on this release and
is stated here so the figures are read correctly.

\subsection{Class selection}

Four classes were retained --- \textbf{Blip}, \textbf{Extremely\_Loud},
\textbf{Scattered\_Light}, \textbf{Violin\_Mode} --- on three criteria
established before any separability analysis, so that the multivariate checks
are a verification rather than the selection rule.

\emph{Abundance.} These are among the most populated of the 24 classes, with
$29\,573$, $24\,002$, $99\,665$ and $2\,392$ members respectively. The smallest
sets the balanced cardinality at $2\,392$ per class, sufficient to train without
augmentation.

\emph{Morphological diversity.} The four span physically distinct families of
instrumental noise~\cite{zevin2017, glanzer2023, bahaadini2018}: short broadband
transients of teardrop morphology and $\sim$0.04\,s duration (Blip);
high-amplitude events that saturate the spectrogram (Extremely\_Loud);
long-duration arches of 2--2.5\,s correlated with ground motion, caused in O3 by
relative motion between the optical suspension and reaction-mass chains
(Scattered\_Light); and narrowband lines at $\sim$500\,Hz and harmonics,
corresponding to resonances of the glass fibres suspending the mirrors
(Violin\_Mode). Spanning the morphological range is what makes the explanation
analysis informative, since each class should elicit a recognizably different
attention pattern.

\emph{Practical constraint.} Four classes keep training and the downstream
quantization grid feasible on a single 6\,GB consumer GPU.

Verification in the tabular domain shows each class isolated by a \emph{different}
descriptor, which is the relevant property: Extremely\_Loud occupies
$\log(\text{SNR}+1)$ between 6 and 12 while the others lie between 2 and 5; Blip
is separated by duration; Scattered\_Light sits below roughly 100\,Hz in peak
frequency with a bimodal bandwidth distribution; and Violin\_Mode is the only
class extending to 2\,kHz. A UMAP projection~\cite{mcinnes2018umap} of the
standardized descriptors confirms the separation independently.

Violin\_Mode is the weakest case, and this is stated because it anticipates a
result that recurs throughout: in the $\log(\text{SNR}+1)$ versus
$\log(\text{duration}+1)$ plane it overlaps both Blip and Scattered\_Light, being
separated only in peak frequency. It is the lowest-recall class at every
configuration tested and the most sensitive class in every explainability
measurement.

Restricting the problem to four classes keeps it tractable, but means the results
are not directly comparable to the full Gravity Spy benchmark.

\subsection{Feature selection}

Of the twelve numerical Omicron descriptors, four were retained on three distinct
criteria.

\emph{Exact redundancy.} The three timestamp variables are mutually correlated at
$r = 1.00$, being the same GPS time to within sub-second offsets. Beyond the
redundancy, retaining any of them would supply \emph{when} a glitch occurred
rather than \emph{what} it is, and the correlation matrix shows this is an
exploitable shortcut: event time correlates at $-0.45$ with the quality factor
and $-0.22$ with peak frequency, reflecting instrumental drift across observing
runs. Separately, central frequency and bandwidth correlate at $0.98$; bandwidth
was retained because peak frequency already encodes spectral position, so
bandwidth supplies extent rather than duplicating location.

\emph{Degenerate distributions.} The nanosecond remainders are uniformly
distributed and correlate with everything between $0.00$ and $0.11$, consistent
with timing quantization noise. The quality factor contains a sentinel value of
order $10^{20}$ that compresses the real distribution to the axis origin.

\emph{Normalized counterpart preferred.} Amplitude is a raw strain amplitude,
heavily skewed and dependent on the calibration and noise level at the time of
the trigger; SNR carries the same information normalized to the local noise floor
and is therefore comparable across detectors and epochs.

The four retained descriptors --- peak frequency, duration, bandwidth, SNR ---
form a complete and non-redundant characterization of a transient in the
time--frequency plane: spectral position, temporal extent, spectral extent and
prominence. Their largest pairwise correlation is $0.39$, well below the $0.98$
that forced the exclusion of central frequency.

Only one of the four image URL columns was kept. The four correspond to
successive time windows of the same detection (0.5\,s, 1\,s, 2\,s, 4\,s), and the
first was retained. Because they are \emph{different views} rather than
duplicates, they are not interchangeable, which is why unavailable images were
excluded rather than substituted from another column.

Residual class imbalance is removed by random undersampling to the cardinality of
the smallest class, giving an exactly balanced set of $9\,568$ images.

\subsection{Preprocessing and partition}

The downloaded spectrograms are matplotlib figures comprising a frame with axes,
colour bar and annotation. A centre crop of $470 \times 550$ pixels isolates the
spectrogram panel, which is then resized to the working resolution; the order
matters, since cropping the original and resizing afterwards keeps the same
portion of the spectrogram framed at every resolution. Channel-wise
normalization uses the standard ImageNet statistics, which also fixes the neutral
reference value used later by the deletion metric.

\textbf{The partition is frozen to disk.} A dedicated script generates the
train/validation/test assignment once, 70\,/\,15\,/\,15 under a fixed seed, and
writes it as a per-file record of path, class and split; every subsequent run
reads that file. A file listed but absent from disk is dropped from its partition
with a warning; a file present but unlisted raises, since an unassigned sample
would otherwise be used without anyone knowing where it belongs.

Freezing the partition rather than recomputing it gives two properties a seed
alone does not. Because the assignment is a function of the file list rather than
of the corpus size, removing a file removes exactly that file instead of
reshuffling every remaining assignment. And the audit trail is externally
verifiable: a reader can confirm that no file appears in two partitions, and the
manifest records the SHA-256 digest of the assignment file.

Four downloaded files proved unreadable. Two were recovered by re-downloading;
two are permanently unavailable at source (0.02\,\% of the corpus) and were
excluded. \textbf{Both belong to the training partition}, so the validation and
test sets are complete and identical to those on which the hyperparameters were
selected, which is why the search did not need to be repeated.

The final corpus is \textbf{9\,566 images, partitioned $6\,695 / 1\,435 /
1\,436$}. Beyond class balance, the three partitions are indistinguishable on the
physical properties of the events: median SNR is $14.97 / 14.09 / 14.98$, with
duration, peak frequency and bandwidth similarly matched.

\section{Model}

\subsection{Architecture}

\texttt{GWGlitchSNN} is a four-block convolutional spiking network. Each block is
a $5 \times 5$ convolution with stride 2 and padding 2 followed by a LIF layer;
spatial downsampling is performed by the stride rather than by pooling, which
keeps the spike tensors dense in time and avoids inserting a non-spiking operator
between layers.

\begin{table}[htbp]
\centering
\begin{tabular}{llll}
\toprule
Stage & Operation & Output shape & Notes \\
\midrule
Input & direct encoding $\times T$ & $[T, B, 3, 112, 112]$ & analog frame repeated \\
Block 1 & Conv $3\to16$, k5, s2, p2 + LIF & $[B, 16, 56, 56]$ & exposed as \texttt{spike1} \\
Block 2 & Conv $16\to32$, k5, s2, p2 + LIF & $[B, 32, 28, 28]$ & exposed as \texttt{spike2} \\
Block 3 & Conv $32\to64$, k5, s2, p2 + LIF & $[B, 64, 14, 14]$ & exposed as \texttt{spike3} \\
Block 4 & Conv $64\to128$, k5, s2, p2 + LIF & $[B, 128, 7, 7]$ & exposed as \texttt{spike4} \\
Readout & Flatten (6\,272) $\to$ Linear $\to$ LIF & $[B, 4]$ & one neuron per class \\
\bottomrule
\end{tabular}
\caption{Layer configuration at the working resolution of 112\,px. The flattened
width is derived from the input size at construction ($25\,088$ at 224\,px,
$6\,272$ at 112\,px, $2\,048$ at 64\,px), and the forward pass raises if the
input resolution does not match the one the model was built for.}
\label{tab:architecture}
\end{table}

The forward pass returns the stacked output spike train \emph{and} the four
hidden spike tensors of shape $[T, B, C, H, W]$, so that every diagnostic and
explainability procedure consumes raw hidden spike trains directly from the model
interface rather than through hooks.

Parameter count is nearly resolution-independent: $295\,332$ at 112\,px, of which
$270\,240$ are convolutional and unaffected by input size. Only the linear
readout shrinks with resolution, a fact with consequences for the design of the
quantization grid discussed in Section~\ref{sec:grid}.

Batch normalization, pooling and dropout are intentionally absent. Normalization
layers interact non-trivially with quantization, and their treatment belongs to a
later phase: batch normalization must be reconsidered for spiking networks in
general, quantizing its parameters is a problem distinct from quantizing weights
and potentials, and two alternative hardware strategies exist --- approximating
it with shifts, or eliminating it entirely.

\subsection{Readout and loss}

Classification uses a \emph{rate-coded readout}: output spikes are summed over
the $T$ time steps and the resulting per-class counts are passed to a
cross-entropy loss as logits. Optimization uses Adam~\cite{kingma2015adam}. This
keeps SAM in the multi-spike regime where its inter-spike-interval formulation
remains meaningful.

\section{Selecting the operating configuration}

Three choices define the operating point: the input resolution, the three
hyperparameters, and the decay constant of the explainer. Each is settled on
evidence and recorded in a single frozen manifest. Taking any of them on test-set
performance would consume the split reserved for the final estimate, so all three
are decided on validation.

\subsection{Input resolution}

Nine training runs were performed: three resolutions (224, 112, 64\,px) at three
seeds each, on the frozen partition, with hyperparameters held fixed and only the
input size varying. Since those hyperparameters had been found at 224\,px, the
protocol systematically handicaps the lower resolutions, and any deficit they
show is an upper bound.

The procedure and its outcome are reported in Section~\ref{sec:resolution}.

\subsection{Hyperparameter search under deployment constraints}

The search uses Optuna~\cite{akiba2019optuna} with a Hyperband
pruner~\cite{li2018hyperband}, maximizing validation accuracy over three
hyperparameters chosen because they govern three distinct axes: learning rate
(optimization stability), $\beta$ (membrane leak and temporal memory), and $T$
(latency, computation and memory cost).

\textbf{Two deployment constraints are imposed inside the search space rather
than by rounding the result.} $\beta$ is restricted to
$\{0.25, 0.5, 0.75, 0.875\}$, values whose decay is realizable as a bit shift,
and the set spans both families ($\beta = 2^{-k}$ and $\beta = 1 - 2^{-k}$)
deliberately. $T$ is capped at 8 for the same reason. A rounded value is one that
was never validated; constraining the search guarantees that whatever it returns
is deployable by construction, and Section~\ref{sec:hyperparameters} shows that
in this case the two differ.

\section{The Spike Activation Map}
\label{sec:sam-method}

SAM is implemented in two stages, following the original
formulation~\cite{kim2021sam}. The temporal spike contribution score weights a
past spike at $t'$ by its recency with respect to the current step $t$:
\begin{equation}
T(t, t') = \exp\left(-\gamma\,|t - t'|\right).
\end{equation}
The neuronal contribution score of a neuron at time $t$ sums that weight over the
neuron's preceding spikes, so a unit that has recently fired repeatedly --- one
with short inter-spike intervals --- receives a high score:
\begin{equation}
N^{k}_{ij,t} = \sum_{t' \in P^{k}_{ij},\; t' < t} T(t, t') .
\label{eq:ncs}
\end{equation}
The heat map at time $t$ masks that score with the spikes actually emitted at $t$
and sums over the channel axis, yielding one spatial map per time step:
\begin{equation}
M_{ij,t} = \sum_{k} N^{k}_{ij,t}\, S^{k}_{ij,t} .
\end{equation}

\paragraph{The summation range is exclusive of $t$.} This is stated explicitly
because an earlier implementation in this work summed over $t' \leq t$, and the
error is instructive. The score measures the contribution of a neuron's
\emph{preceding} trajectory; including $t' = t$ adds a term of weight
$\exp(0) = 1$ for every neuron firing at $t$, so the map produced is the correct
SAM \emph{plus the raw binary spike map}. The pedestal is independent of both
$\gamma$ and the spike history, and therefore inflates every correlation between
maps computed at different $\gamma$. A consequence of the correct definition is
that SAM at $t = 0$ is identically zero, no past history existing at the first
step.

The implementation uses the recursion
\begin{equation}
N_{t+1} = e^{-\gamma}\left(N_{t} + S_{t}\right), \qquad N_0 = 0,
\end{equation}
algebraically identical to Equation~\ref{eq:ncs} but $O(T)$ rather than
$O(T^2)$. A reference implementation computing the explicit sum is retained and
the two are asserted equal at run time.

Visualizations overlay every time step on the denormalized input under a
\emph{shared colour scale}, so that a change in intensity across time is genuine
and not an artefact of per-panel rescaling.
